\documentclass{article}

\usepackage{tabularx}
\usepackage{booktabs}

\title{Problem Statement and Goals\\\progname}

\author{\authname}

\date{}

\input{../Comments.text}
\input{../Common.text}

\begin{document}

\maketitle

\begin{table}[hp]
\caption{Revision History} \label{TblRevisionHistory}
\begin{tabularx}{\textwidth}{llX}
\toprule
\textbf{Date} & \textbf{Developer(s)} & \textbf{Change}\\
\midrule
Date1 & Name(s) & Description of changes\\
Date2 & Name(s) & Description of changes\\
... & ... & ...\\
\bottomrule
\end{tabularx}
\end{table}

\section{Problem Statement}

\wss{You should check your problem statement with the
\href{https://smiths.github.io/capTemplate/Checklists/ProbState-Checklist.pdf}
{problem statement checklist}}. 

\wss{You can change the section headings, as long as you include the required
information.}

\subsection{Problem}

\wss{What problem are you trying to solve?  Why is it important? Do not talk
about how you will solve the problem, only what the problem is.  AI is rarely
part of the problem.  The problem is ``what'' you want to do, not ``how'' you
want to do it.  The problem statement should be written in a way that is
understandable to a non-technical audience.}

\subsection{Inputs and Outputs}

\wss{Characterize the problem in terms of ``high level'' inputs and outputs.  
Use abstraction so that you can avoid details.}

\subsection{Stakeholders}

\subsection{Environment}

\wss{Hardware and Software Environment for the users.  The developer environment
is summarized as part of the developer plan.}

\section{Goals}

\section{Stretch Goals}

\wss{Goals you hope to get to, but not necessary.  They are also helpful if the
scope of the original project needs to be expanded.}

\section{Custom Documents (CDs)}

\wss{For CAS 741: State whether the project is a research project. This
designation, with the approval (or request) of the instructor, can be modified
over the course of the term.}

\wss{For SE Capstone: List your custom documents (CDs).  Potential CDs include
usability testing, code walkthroughs, user documentation, formal proof,
GenderMag personas, Design Thinking, etc.  (The full list is on the course
outline and in Lecture 02.) Normally the number of CDs will be two (one per
term).  Approval of the CDs will be part of the discussion with the instructor
for approving the project. The CDs, with the approval (or request) of the
instructor, can be modified over the course of the term.}

\newpage{}

\section{Appendix --- Reflection}

\subsection*{Jake Finlay}

\subsubsection*{What went well while writing this deliverable?}

What made this deliverable straightforward for me was the amount of discussion
we had already had with our TA and our stakeholders before writing anything
down. By the time we sat down to write, the problem was well understood by
everyone. Gathering the information for the deliverable was mostly a matter
of writing down what we had already agreed on rather than working it out from
scratch. The other thing that helped was that we had broken the deliverable into
issues on GitHub ahead of time and assigned them to team members. That meant
there was never any ambiguity about who owned what or what needed to be picked
up next. This allowed me to move straight into writing instead of spending time
coordinating.

\subsubsection*{What pain points did you experience during this deliverable,
and how did you resolve them?}

The main pain point for me was collaborating through GitHub. Our team had not
yet agreed on conventions for pull requests and commit messages, so everyone
took a different approach. This made reviews take much longer than they needed
to. The same issue appeared in the deliverables themselves. Each person's
writing was strong on its own, but details such as formatting varied between
team members and had to be aligned. We resolved this by discussing these
differences in pull request comments and in our team chat. These discussions
were necessary because our development plan was still being written and we had
nothing to follow at the time of writing.

\end{document}